
\documentclass[11pt,a4paper]{article}
\usepackage[margin=2.5cm]{geometry}
\usepackage[utf8]{inputenc}
\usepackage[T1]{fontenc}
\usepackage[french]{babel}
\usepackage{enumitem}
\usepackage{hyperref}


\begin{document}
	\pagestyle{empty}
	
	\noindent
	Friedly WOLI \\  % \hfill Airbus \\
	Toulouse, 31100  \\
	0767339290 \\
	friedlysedjro@gmail.com 
	
	\begin{flushright}
		
		Université de Toulouse \\
		118 route de Narbonne 31062 
	\end{flushright}
	
	
	\bigskip
	
	
	
	
	
	
	\textbf{Objet :} Candidature master MSME \\  \\
Madame, Monsieur, \\ 

Actuellement étudiant en mécanique énergétique à l’Université de Toulouse, je souhaite intégrer le Master Modélisation et Simulation en Mécanique et Énergétique (MSME) en alternance afin d’approfondir ma formation en mécanique théorique et appliquée tout en développant une expérience concrète en entreprise dans le domaine de la modélisation et de la simulation.

Au cours de mon parcours universitaire, j’ai acquis de solides bases en mécanique des fluides, en mécanique des solides et des structures ainsi qu’en transferts thermiques. J’ai également développé un intérêt particulier pour la formulation mathématique des modèles physiques et leur analyse. Les enseignements en mathématiques appliquées et en calcul scientifique m’ont permis de mieux appréhender la rigueur nécessaire à la modélisation et à la simulation de systèmes physiques complexes.

Le parcours MSME en alternance correspond pleinement à mes objectifs. L’équilibre entre formation académique avancée et immersion en entreprise représente pour moi une opportunité idéale pour consolider mes connaissances théoriques tout en les appliquant à des problématiques industrielles concrètes. Cette articulation entre théorie et pratique constitue un cadre particulièrement formateur pour développer des compétences solides en modélisation physique et en simulation numérique.

Mon objectif professionnel est de travailler dans le domaine de la simulation et de la modélisation appliquées aux systèmes mécaniques et énergétiques, notamment dans les secteurs industriels à forte intensité technologique comme l’aéronautique. L’alternance me permettra de développer une compréhension plus concrète des enjeux industriels tout en renforçant mes compétences scientifiques et techniques.

Rigoureux, curieux et persévérant, je suis prêt à m’investir pleinement dans cette formation exigeante et dans les missions qui me seront confiées en entreprise. Intégrer le Master MSME en alternance représente pour moi une étape essentielle pour construire un profil alliant solide formation scientifique et expérience professionnelle en simulation numérique. 

Je vous remercie par avance de l’attention que vous porterez à ma candidature et vous prie d’agréer, Madame, Monsieur, l’expression de mes salutations distinguées.



	

	
	
	
	\vspace{0.8cm}
	
	
	
	
	
	
	\begin{flushright}
		
		\textbf{Friedly WOLI}
	\end{flushright}
	
	
\end{document}




